\documentclass[11pt,a4paper]{article}
\usepackage[T1]{fontenc}
\usepackage{lmodern,amsmath,amssymb,amsthm,booktabs,microtype}
\usepackage[margin=27mm]{geometry}
\usepackage[colorlinks=true,linkcolor=blue,citecolor=blue,urlcolor=blue]{hyperref}
\setlength{\emergencystretch}{2em}
\newtheorem{theorem}{Theorem}[section]
\newtheorem{lemma}[theorem]{Lemma}
\newtheorem{proposition}[theorem]{Proposition}
\newcommand{\Q}{\mathbb Q}\newcommand{\Z}{\mathbb Z}
\DeclareMathOperator{\Tr}{Tr}
\title{Five local factors with squared reciprocal tests}
\author{}\date{17 November 2026}
\begin{document}\maketitle

\begin{abstract}
Five explicit factors separate the valuations of the comparison and
reciprocal-abscissa tests at every negative odd valuation of the
contact weight. Coupled to the retained trace and pole equations,
they recover the original comparison depths. We prove the selected
contact residues, fixed-place limits, real positivity and nonvanishing
on every prepared real character. These are arithmetic interface
lemmas. No quadratic compilation, coordinate reduction, or independent
completeness theorem for a new polynomial is asserted.
\end{abstract}

\section{Tests, factors and inherited equations}

Let $K=\Q(\sqrt2)$, put $D=11-6\sqrt2$, and use
$E:y^2=x^3-Dx$. For a point with finite nonzero abscissa set
$f=1/x$, and set $f(O)=0$. The physical tests of
\cite{General,Charts,Explicit} are
\[
 D_H=f(H),\quad A_i=f(U_i),\quad
 B_i=\overline{f(V_i)}\ (1\le i\le4),\quad
 B_0=\overline{f(H-P_0)}.
\]
Keep the original cross comparisons $C_1,\ldots,C_4$, input
comparison $I$, and affine Lucas residual
\[
 G=L_3^2-(L_2+1)(L_2-1)L_3L_4+(L_2-1)^2L_4^2-1.
\]
All these tests are elements of $K$ on the retained point and
quotient domain. Zero tests are permitted.

Let $q\in\Q\setminus\{0,-1\}$ and $R=(1+q)^{-1}$.
Define
\begin{equation}\label{eq:factors}
\begin{aligned}
 \Phi_0&=1+q^{11}C_2^4+q^2D_H^2+qB_1^2,\\
 \Phi_1&=1+q^{11}G^4+q^2A_3^2+qB_0^2,\\
 \Phi_2&=1+q^4I^8+qA_2^2+q^2A_4^2+q^3B_3^2,\\
 \Phi_3&=1+q^{11}C_3^4+q^2C_1^4+qB_4^2,\\
 \Phi_4&=1+q^{11}C_4^4+q^2A_1^2+qB_2^2.
\end{aligned}
\end{equation}
Write $W=\prod_{i=0}^4\Phi_i=A+B\sqrt2$. The inherited
rational trace and pole-input equations are
\begin{equation}\label{eq:trace}
 A^2-2A\eta-2B^2=0,\qquad (1+q)b=\kappa^2\eta^2,
                       \qquad \kappa=408/409.
\end{equation}
Retain the source pole implication that $v_p(b)$ is not a negative
odd integer at a consumed prime. This paper proves the consequences
of that implication; it does not replace the pole predicate or its
realization theorem.

\section{Arbitrary odd contact depth}

Fix an odd rational prime $p$ unramified in $K$, with $D$ a unit
at its places and $v_p(\kappa)\le0$. These conditions hold at the
consumed primes of the retained sources. Normalize every place by
$v(p)=1$, so valuations of elements of $K$ are integral.
Suppose
\[
                         v_p(q)=-\mu<0,\qquad \mu\text{ odd}.
\]
Then $v_p(R)=\mu$. Assign valuation infinity to zero tests and
omit their terms in every comparison of finite valuations.

\begin{lemma}\label{lem:even}
Every finite reciprocal-abscissa valuation at a place above $p$
is even.
\end{lemma}
\begin{proof}
For an affine point put $h=y/x^2$. Then $h^2=f-Df^3$.
If $v(f)>0$, the right side has valuation $v(f)$; if $v(f)<0$,
it has valuation $3v(f)$. Either must be even. The case $v(f)=0$
is immediate. The identity convention $f=0$ gives infinity.
\end{proof}

\begin{lemma}\label{lem:separate}
At each place above $p$, every factor in \eqref{eq:factors} is
nonzero and has valuation at most zero. Its minimum term is unique.
Its inverse is integral and has residue zero or one. More precisely,
a negative factor valuation gives inverse residue zero; a zero factor
valuation gives factor and inverse residue one.
\end{lemma}
\begin{proof}
For each of $\Phi_0,\Phi_1,\Phi_3,\Phi_4$, the finite term
valuations lie in the distinct classes
\[
                  0,\quad\mu,\quad2,\quad-\mu\pmod4.
\]
Here the third term uses an even reciprocal valuation, or the fourth
power $C_1^4$. The classes are distinct because $\mu$ is odd.

For $\Phi_2$, the valuation classes modulo eight lie in the
following disjoint sets:
\[
\begin{array}{c|c}
\text{term}&\text{possible classes modulo }8\\\hline
1&\{0\}\\
q^4I^8&\{4\}\\
qA_2^2&\{-\mu,-\mu+4\}\\
q^2A_4^2&\{-2\mu,-2\mu+4\}\\
q^3B_3^2&\{-3\mu,-3\mu+4\}.
\end{array}
\]
They partition the eight classes. Omitting zero terms preserves
separation. The constant term ensures that the unique minimum is
at most zero. If the minimum is zero, every other term has strictly
positive valuation, so the factor is one modulo the place.
\end{proof}

\begin{theorem}\label{thm:extract}
Assume the rational equations \eqref{eq:trace} and the inherited
pole implication at $p$. Then at some place above $p$,
\[
 \begin{gathered}
 v(C_2,C_3,C_4,G)>21\mu/8,\qquad
 v(C_1)>3\mu/8,\qquad v(I)>7\mu/16,\\
 v(D_H,A_1,A_3,A_4)>3\mu/4,\\
 v(A_2,B_0,B_1,B_2,B_4)>\mu/4,\qquad
 v(B_3)>5\mu/4.
 \end{gathered}
\]
In particular the source thresholds hold:
\[
 v(C_2,C_3,C_4,G)\ge3,\qquad v(C_1,I)\ge1,
 \qquad v(D_H,A_i,B_i,B_0)\ge2.
\]
No nonzero condition on an individual test is required.
\end{theorem}
\begin{proof}
Lemma~\ref{lem:separate} gives $W\ne0$ and
$A^2-2B^2\ne0$. The first trace equation forces $A\eta\ne0$.
Thus, without imposing a new global guard,
\[
 V:=\Tr(W^{-1})=\frac{2A}{A^2-2B^2}\ne0,\qquad
                   \eta=V^{-1},\qquad b=\frac{\kappa^2R}{V^2}.
\]
Put $\ell=\max_{v\mid p}v(W)\le0$. Then $v_p(V)\ge-\ell$,
so $v_p(b)\le\mu+2\ell$. This valuation is odd, and the pole
implication makes it positive. Consequently $\ell>-\mu/2$.
At a place attaining $\ell$, each factor has valuation at least
$\ell$, because the remaining factor valuations are nonpositive.
Each finite summand therefore has valuation greater than $-\mu/2$.

For a term $q^eT^m$ this is precisely
\[
                             v(T)>\frac{e-1/2}{m}\mu.
\]
The relevant coefficients are
\[
\begin{array}{c|cccccc}
\text{term}&q^{11}C^4&q^2C_1^4&q^4I^8&qT^2&q^2T^2&q^3T^2\\\hline
\text{coefficient}&21/8&3/8&7/16&1/4&3/4&5/4.
\end{array}
\]
This proves the strict bounds. Integrality gives the stated comparison
thresholds, and Lemma~\ref{lem:even} gives reciprocal threshold two.
An infinite valuation satisfies every bound. The possible negative
valuation of $\kappa$ at 409 only strengthens the trace inequality.
\end{proof}

The theorem supplies exactly the depths used by the retained
formal-comparison, Lucas and input-extraction arguments. Its use
requires those arguments and their original exceptional sets; it
introduces no new exceptional prime or radical-height loss.
When a detector has no consumed prime, this local argument imposes
no additional nonzero guard elsewhere.

\section{Selected contacts and preparation limits}

\begin{proposition}\label{prop:selected}
Suppose $v_p(q)=-1$, $v_p(\kappa)=0$, and at one place above $p$
all comparisons $C_i,I,G$ have valuation at least four and every
reciprocal test has valuation at least two. Then $W$ is one modulo
that place. At the opposite place $W^{-1}$ has residue zero or one.
The actual reciprocal trace $V$ is a unit, of residue one or two,
and the input in \eqref{eq:trace} has $v_p(b)=1$.
\end{proposition}
\begin{proof}
In the displayed order of the nonconstant terms, their valuation
lower bounds are
\[
\begin{array}{c|c}
\Phi_0&5,2,3\\
\Phi_1&5,2,3\\
\Phi_2&28,3,2,1\\
\Phi_3&5,14,3\\
\Phi_4&5,2,3.
\end{array}
\]
All are positive, including the cases where a term is zero.
The opposite-place assertion follows from Lemma~\ref{lem:separate}
without any opposite-place formalness assumption. At a split prime,
the two inverse residues sum to one or two. At an inert prime,
conjugation preserves the unique place and fixes successful residue
one, so the trace residue is two. The prime is odd, so the trace
is a unit. Since $v_p(R)=1$, the assertion about $b$ follows.
\end{proof}

\begin{proposition}\label{prop:preparation}
The following preparation statements hold for the physical factors.
\begin{enumerate}
\item On any fixed finite-place neighborhoods where all tests are
regular and bounded, $q\to0$ gives $W\to1$, $V\to2$ and
$b\to(204/409)^2$. The convergence is uniform on those neighborhoods.
\item At every prepared real character where the tests are finite
and $q>0$, all five factors are at least one, $W>0$ and $V>0$.
For bounded real tests and $q\to0^+$, one has
$b\to(204/409)^2\in(0,17)$.
\item Suppose the unchanged contact preparation makes the only odd
valuations of $R$ occur at the selected primes and supplies the
hypotheses of Proposition~\ref{prop:selected}. For the resulting
actual nonzero traces, $b$ has square reduced denominator.
\end{enumerate}
\end{proposition}
\begin{proof}
Each nonconstant term in \eqref{eq:factors} has a positive power
of $q$. A common lower valuation bound for its powered tests shows
that sufficiently deep $q\to0$ makes every term arbitrarily small.
Zero tests cause no restriction. The assertion is valid at ramified
fixed places as well, since no residue-class separation is used there.
This gives $W\to1$ and $V\to2$. Since $R=(1+q)^{-1}\to1$,
the limit of $b$ is the stated rational square.

Its valuations at 2, 3 and 17 are respectively $4,2,2$.
Taking the preparation sufficiently deep preserves these depths
and the required local square classes, since
\[
                  \frac{b}{(204/409)^2}=R(2/V)^2
\]
and $R$ can lie in the local square neighborhood of one.
Thus the retained finite-place pole preparation is compatible with
the changed factors. The point family and single source slope are
unchanged; no independently selected trace is used.

At a real embedding every test exponent is even. For $q>0$, each
term is nonnegative and each factor is at least one. This proves
nonvanishing separately at every prepared real character, including
both embeddings of $K$, and gives a positive reciprocal trace.
The limits are uniform on bounded real neighborhoods. In particular
a sufficiently large retained slope cutoff gives $0<b<17$.
This does not assert nonvanishing at arbitrary nonreal characters.

Finally Proposition~\ref{prop:selected} gives valuation one at
each selected prime. At every other prime the assumed contact
factorization gives even valuation of $R$. Multiplication by
$\kappa^2$ and division by $V^2$ preserve parity. All negative
valuations of $b$ are therefore even.
\end{proof}

The inherited ordered preparation first fixes the index labels and
bounded neighborhoods, then the contact primes and final reference
progression, and only then realizes the pole witnesses for the
actual inputs. The preceding propositions preserve its local and
real requirements after increasing the fixed precision and slope
cutoffs. They do not prove that the new factors fit a quadratic
table, or that an unverified new chart preserves every source
hypothesis. No complete defining polynomial or coordinate reduction
follows from the physical-factor lemmas alone.

\section{Alternative pairings and exact checks}

The same valuation proof permits two further pairings, useful when
distributing test outputs. The terms $q^9C^4$ and $q^{11}C'^4$
have opposite odd classes modulo four. A third term $q^2T$, with
$v(T)$ divisible by four, occupies class two. The weaker comparison
bound is $v(C)>17\mu/8$, still at least three; its prepared
valuation lower bound is seven. Likewise $qT$ and $q^3T'$ have
opposite odd classes whenever $v(T),v(T')$ are divisible by four.
Such outputs include reciprocal squares and $C_1^4$.
The special pattern $q^4I^8$ together with $qT_1,q^2T_2,q^3T_3$
has the same modulo-eight partition as $\Phi_2$.
These are arithmetic statements, not quadratic-feasibility claims.

The exact standard-library checker is
\nolinkurl{computation/quartic/local_detector_factors_v66/check.py}.
It verifies every residue pattern over the complete period
$\mu\bmod8$, all strict rational coefficients, the sixteen
selected-contact lower bounds, and the fixed-place limit depths.
The unbounded valuation, trace and preparation arguments are proved
above and are not inferred from sampled valuations.

\begin{thebibliography}{9}\raggedright
\bibitem{General}\emph{A shifted Lucas equation and a quartic in 200 variables},
repository manuscript, version 33, local extraction and completeness,
\nolinkurl{papers/quartic_existence/quartic200_20261115_v33.tex}.
\bibitem{Explicit}\emph{Three-detector quartics with 229 rational coordinates},
repository manuscript, version 62, complete detector comparison and
complete witnesses,
\nolinkurl{papers/quartic/quartic229_20261117_v62.tex}.
\bibitem{Charts}\emph{A quartic in 204 variables from recovered elliptic differences},
repository manuscript, version 27,
\nolinkurl{papers/quartic_existence/quartic204_20261114_v27.tex}.
\end{thebibliography}
\end{document}
